\documentclass[aspectratio=169,12pt]{beamer}
\usepackage{amsmath}
\usepackage{amssymb}
\usepackage{array}
\usepackage[most]{tcolorbox}
\usepackage{graphicx}
\usepackage[sfdefault,lf]{FiraSans}
\renewcommand{\familydefault}{\sfdefault}
\setlength{\parindent}{0pt}
\everymath{\displaystyle}

\setbeamertemplate{navigation symbols}{}
\setbeamertemplate{footline}{}
\setbeamertemplate{headline}{}
\setbeamertemplate{blocks}[rounded][shadow=false]

\definecolor{mmpageWhite}{HTML}{FCFBF7}
\definecolor{mmtanSoft}{HTML}{F7F5EF}
\definecolor{mmgreenDeep}{HTML}{244B2C}
\definecolor{mmgreenLeaf}{HTML}{5D8A56}
\definecolor{mmgreenSoft}{HTML}{E4EFD9}
\definecolor{mmtext}{HTML}{163221}
\definecolor{mmwarn}{HTML}{8F2F36}
\definecolor{mmwarnSoft}{HTML}{F8E8EA}

\setbeamercolor{background canvas}{bg=mmpageWhite}
\setbeamercolor{normal text}{fg=mmtext,bg=mmpageWhite}
\setbeamercolor{itemize item}{fg=mmgreenDeep}
\setbeamercolor{itemize subitem}{fg=mmgreenDeep}

\newtcolorbox{MMCard}[2][]{enhanced,breakable=false,colback=#2,colframe=black,boxrule=1.5pt,arc=8pt,left=5pt,right=5pt,top=4pt,bottom=4pt,before skip=0pt,after skip=0pt,#1}
\newcommand{\KoruIcon}{\raisebox{-0.2em}{\includegraphics[height=1.05em]{../../../manamaths/SITE/header-logo.png}}}
\newcommand{\NotesTitle}[1]{{\colorbox{mmtanSoft}{\parbox{0.975\linewidth}{\vspace{0.14em}\hspace{0.25em}{\LARGE\bfseries\textcolor{mmgreenDeep}{#1}\hfill\KoruIcon}\vspace{0.14em}}}\par\vspace{0.18em}{\color{mmgreenLeaf}\rule{\linewidth}{2.0pt}}\par\vspace{0.12em}}}
\newcommand{\SectionTitle}[1]{{\large\bfseries\textcolor{mmgreenDeep}{#1}}\par\vspace{0.04em}}
\newcommand{\GreenDivider}{\par\vspace{0.01em}{\color{mmgreenLeaf}\rule{\linewidth}{2.4pt}}\par\vspace{0.03em}}
\newcommand{\KeyIdea}[1]{\begin{MMCard}[title={\textbf{Key idea}},colbacktitle=mmgreenSoft,coltitle=mmgreenDeep]{mmtanSoft}\small #1\end{MMCard}}
\newcommand{\WorkedExample}[2]{\begin{MMCard}[title={\textbf{#1}},colbacktitle=mmgreenSoft,coltitle=mmgreenDeep]{white}\small #2\end{MMCard}}
\newcommand{\CommonMistake}[1]{\begin{MMCard}[title={\textbf{Common mistake}},colbacktitle=mmwarnSoft,coltitle=mmwarn]{white}\small #1\end{MMCard}}
\newcommand{\TryThese}[1]{\begin{MMCard}[title={\textbf{Try these}},colbacktitle=mmgreenSoft,coltitle=mmgreenDeep]{mmtanSoft}\small #1\end{MMCard}}
\newcommand{\StepLine}[2]{\textbf{#1.} #2\par\vspace{0.03em}}

\usepackage{tikz}
\setbeamertemplate{background}{
 \begin{tikzpicture}[remember picture,overlay]
 \draw[black,line width=1.5pt,rounded corners=10pt]
 ([xshift=0.45cm,yshift=-0.45cm]current page.north west)
 rectangle
 ([xshift=-0.45cm,yshift=0.45cm]current page.south east);
 \end{tikzpicture}
}

\begin{document}

\begin{frame}[t,shrink=10]
\NotesTitle{Notes \& Steps}
\footnotesize
\begin{columns}[T,onlytextwidth]
\begin{column}{0.48\textwidth}
\KeyIdea{Capacity measures how much a container can hold. The common units are millilitres (mL), litres (L), and kilolitres (kL). 1000 mL = 1 L, 1000 L = 1 kL, and 1 mL of water weighs 1 g.}
\SectionTitle{Steps}
\StepLine{1}{Identify the unit you have and the unit you need.}
\StepLine{2}{To convert L \(\to\) mL, multiply by 1000.}
\StepLine{3}{To convert mL \(\to\) L, divide by 1000.}
\StepLine{4}{For kL, multiply or divide by 1000 again.}
\StepLine{5}{Check: does your answer make sense? A bathtub is about 150 L, not 150 mL.}
\end{column}
\begin{column}{0.48\textwidth}
\SectionTitle{Capacity benchmarks}
\begin{itemize}
  \item \textbf{1 mL}: a drop of water, a teaspoon
  \item \textbf{5 mL}: a standard teaspoon
  \item \textbf{250 mL}: a standard cup
  \item \textbf{330 mL}: a can of drink
  \item \textbf{500 mL}: a sports bottle
  \item \textbf{1 L}: a carton of milk
  \item \textbf{1.5--2 L}: a water bottle
  \item \textbf{5--10 L}: a bucket
  \item \textbf{100--200 L}: a bath
  \item \textbf{1 kL}: a small water tank
  \item \textbf{10 kL}: a swimming pool (small)
\end{itemize}
\end{column}
\end{columns}
\GreenDivider
\begin{columns}[b,onlytextwidth]
\begin{column}{0.48\textwidth}
\CommonMistake{Confusing mL and L. A can of drink is 330 mL, not 330 L (that would fill several bathtubs). Always imagine the container.}
\end{column}
\begin{column}{0.48\textwidth}
\TryThese{\begin{enumerate}
  \item Convert 2.5 L to mL.
  \item Convert 1500 mL to L.
  \item How many 250 mL cups in a 2 L bottle?
\end{enumerate}}
\end{column}
\end{columns}
\end{frame}

\begin{frame}[t,shrink=12]
\NotesTitle{Notes \& Steps}
\begin{columns}[T,onlytextwidth]
\begin{column}{0.48\textwidth}
\WorkedExample{Example 1: L to mL}{A bottle holds 1.75 L of water. Convert to mL.\[1.75 \times 1000 = 1750\text{ mL}\]}
\end{column}
\begin{column}{0.48\textwidth}
\WorkedExample{Example 2: mL to L}{A jug contains 850 mL of juice. How many litres?\[850 \div 1000 = 0.85\text{ L}\]}
\end{column}
\end{columns}
\GreenDivider
\begin{columns}[T,onlytextwidth]
\begin{column}{0.48\textwidth}
\WorkedExample{Example 3: filling glasses}{A 2 L bottle fills 250 mL glasses. How many?\[2000 \div 250 = 8\text{ glasses}\]}
\end{column}
\begin{column}{0.48\textwidth}
\WorkedExample{Example 4: kL conversion}{A water tank holds 3.5 kL. How many litres?\[3.5 \times 1000 = 3500\text{ L}\]}
\end{column}
\end{columns}
\end{frame}

\end{document}
